\subsection*{The actual quartic origin at class two and its complete toric image}
\begin{statement}{E.105}{Proved}{A class-two lift of the quartic coset origin, with an explicit continuous toric criterion}
Retain the plus quartic, its twelve cusps, the field
$F=\mathbb Q(\zeta _8)$, and every actual coefficient map of E.103.
Let $R$ be $\mathbb Z_\ell$ for an odd prime, the product of these
rings over all odd primes, or $\mathbb Q_\ell$ for any prime.
Write $A=A_R$ and $W=T_RT_D\subset A$.
The geometric class-two group and its central kernel are
\[
 N_{R,2},\qquad B_R=\gamma _2N_R/\gamma _3N_R
       =\Lambda ^2_R A_R,\qquad
 \operatorname{rank}A_R=17,\quad\operatorname{rank}B_R=136.
 \tag{E105.1}
\]
For integral coefficients this is the indicated geometric
pro-$\ell$, or product of pro-$\ell$, quotient of $\pi _1(U)$.
For rational coefficients it is the geometric two-step Malcev
pushout specified in the proof.  Let $\mathcal E_{R,2}$ be the
resulting actual arithmetic extension.

The section $s_R$ constructed in E.103 lifts to a continuous
section $S_R:G_{\mathbb Q}\to\mathcal E_{R,2}$.
This statement applies to the actual choices of coset lifts and
unit roots in E.103: (E105.18) transports the construction to
each such choice.  No new abelian origin is substituted.

Let $\iota:R(1)^2\to A_R$ be the actual coefficient section in
(E103.16), and write $\alpha,\beta$ for continuous Kummer classes.
The complete image of the class-two section space under the
coordinate units $(x/z,y/z)$ is
\[
 \left\{(\alpha,\beta)\in H^1_{\rm cts}(\mathbb Q,R(1))^2:
             \alpha\cup\beta=0
                  \text{ in }H^2_{\rm cts}(\mathbb Q,R(2))\right\}.
 \tag{E105.2}
\]
Every pair in this set is attained above the specific E.103
abelian section $\iota(\alpha,\beta)s_R$.
The following formulas evaluate the condition in (E105.2) on
every completed input, including unrestricted ramification.

First let $R=\mathbb Z_\ell$, with $\ell$ odd.  Write
\[
 \alpha=\sum_p a_p\kappa_p,\qquad
 \beta=\sum_p b_p\kappa_p,
 \quad a_p,b_p\in\mathbb Z_\ell,\quad a_p,b_p\longrightarrow0,
 \tag{E105.3}
\]
where $p$ runs over the rational primes and the sums converge
uniformly as cocycles.  For $q\ne\ell$ put $e_q=v_\ell(q-1)$.
Choose a primitive $\ell^{e_q}$th root $\eta_q\in\mathbb F_q^*$,
and, for $p\ne q$, define $r_{q,p}\in\{0,\ldots,\ell^{e_q}-1\}$ by
\[
 \overline p^{(q-1)/\ell^{e_q}}=\eta_q^{r_{q,p}}.
 \quad
 A_q=\sum_{p\ne q}a_pr_{q,p},\quad
 B_q=\sum_{p\ne q}b_pr_{q,p},\quad
 c_q=-\frac{A_qb_q-B_qa_q}{q-1}.
 \tag{E105.4}
\]
When $e_q=0$, all $r_{q,p}$ are zero.
Then (E105.2) is exactly the condition
\[
          c_q\in\mathbb Z_\ell\quad(q\ne\ell),
                  \qquad c_q\longrightarrow0.
 \tag{E105.5}
\]
For the product of all odd coefficient rings apply (E105.5)
in every component, with the uniform congruence topology of
(E103.5).  For $R=\mathbb Q_\ell$, including $\ell=2$, use
$a_p,b_p\in\mathbb Q_\ell$ tending to zero and the same formula
(E105.4); the exact condition is
\[
                         c_q\longrightarrow0
                               \text{ in }\mathbb Q_\ell.
 \tag{E105.6}
\]
This includes a common bound on the denominators of the $c_q$.
The conditions are independent of the specified roots and
residue generators through the explicit changes in the proof.

For rational numbers $a,b\in\mathbb Q^*$ and odd integral
coefficients, (E105.5) is the finite test that, at every prime $q$,
the $\ell$-primary component of
\[
 (-1)^{v_q(a)v_q(b)}
       \frac{a^{v_q(b)}}{b^{v_q(a)}}\pmod q\quad\text{in }\mathbb F_q^*
 \tag{E105.7}
\]
is trivial.  Only the union of the numerator and denominator
prime supports of $a$ and $b$ needs to be checked.
For the product of all odd coefficients the order of every
residue in (E105.7) must be a power of two.
Every rational-number pair passes the rational-coefficient
test (E105.6).  In contrast, (E105.29) below is an actual
completed pair that passes every finite-coefficient and every
local test and fails (E105.5) and (E105.6).
The pair $(a,b)=(2,7)$ fails the integral test at $\ell=3,q=7$.

Over each specific liftable abelian section, the entire
geometric-conjugacy fibre is the torsor under
$H^1_{\rm cts}(\mathbb Q,B_R)$ with the literal action
$(z,S)\mapsto(g\mapsto z(g)S(g))$.
Its coefficient maps include the following actual direct sum:
\[
 \Lambda^2 A_R
  =\Lambda^2W\ \oplus\ (W\otimes(T_RE_+)^3)
       \ \oplus\ R(1)^6\ \oplus\bigl(\operatorname{Sym}^2T_RE_+\bigr)^3.
 \tag{E105.8}
\]
The ranks are $55,66,6,9$.
The theorem evaluates the origin and the complete toric image;
it retains the further obstruction for the other abelian
preimages in (E103.3), and the unevaluated cohomology factors
in (E105.8).  The full arbitrary-input quartic class-two
classification in E.23 therefore remains active.
The source map to the integral two-primary abelian quotient
still carries E.18's obstruction; no section of the full
profinite quartic extension is inferred.
\end{statement}
\begin{proof}
\emph{Literal cusps, frames, and fourth-division paths.}
Put $\zeta=\zeta_8$, $\alpha_i=\zeta^{1+2i}$, and use
\[
 \begin{array}{c|c|c|c}
 d&\text{point}&s_d&L_d\\ \hline
 A_i&[0:\alpha_i:1]&x/z&y-\alpha_i z\\
 B_i&[\alpha_i:0:1]&y/z&x-\alpha_i z\\
 C_i&[\alpha_i:1:0]&z/y&x-\alpha_i y
 \end{array}
 \qquad O=B_0.
 \tag{E105.9}
\]
All tangents have value one.  In each displayed affine chart
the nonzero coordinate is $\alpha_i(1+s_d^4)^{1/4}$, where
the fourth root has constant coefficient one.  Consequently
the leading value of $L_d$ at $d$ is $\alpha_i/4$, in order
four.  Substituting $L_d=0$ in the quartic equation leaves
the fourth power of the omitted coordinate, proving
$\operatorname{div}(L_d/L_O)=4d-4O$ including multiplicities.

Let $z_d$ be the framed point $\mathcal O(d-O)$ of $J_{C,D}$,
with vector frames $s_d^{-1}$ at $d$, $s_O$ at $O$, and the
canonical section at the other cusps.  Write $f_d=L_d/L_O$.
Multiplication by $f_d$ trivializes $4z_d$ and gives the
actual torus coordinates
\[
 t_{q,d}=
 \frac{\operatorname{lead}_q(f_d)}{\operatorname{lead}_O(f_d)},
       \qquad q,d\ne O.
 \tag{E105.10}
\]
Here leading values use the exact parameters in (E105.9),
with the appropriate zero or pole order removed.
The program evaluates all $121$ entries in
$\mathbb Q[\zeta]/(\zeta^4+1)$ and retains their four rational
coefficients and their exact root phases:
$t_{q,d}=\zeta^{K_{qd}}(\sqrt2)^{N_{qd}}$.
In the order
\[
        A_0,A_1,A_2,A_3,B_1,B_2,B_3,C_0,C_1,C_2,C_3
 \tag{E105.11}
\]
the integer matrix $N$ is
\[
\left(\begin{array}{rrrrrrrrrrr}
-8&-3&-2&-3&-5&-6&-5&-4&-4&-4&-4\\
-3&-8&-3&-2&-5&-6&-5&-4&-4&-4&-4\\
-2&-3&-8&-3&-5&-6&-5&-4&-4&-4&-4\\
-3&-2&-3&-8&-5&-6&-5&-4&-4&-4&-4\\
-5&-5&-5&-5&-10&-6&-4&-5&-5&-5&-5\\
-6&-6&-6&-6&-6&-12&-6&-6&-6&-6&-6\\
-5&-5&-5&-5&-4&-6&-10&-5&-5&-5&-5\\
-4&-4&-4&-4&-5&-6&-5&-8&-3&-2&-3\\
-4&-4&-4&-4&-5&-6&-5&-3&-8&-3&-2\\
-4&-4&-4&-4&-5&-6&-5&-2&-3&-8&-3\\
-4&-4&-4&-4&-5&-6&-5&-3&-2&-3&-8
\end{array}\right).
 \tag{E105.12}
\]
The retained $K_{qd}\in\{0,\ldots,7\}$, rather than a
discarded root-of-unity factor, is used below for $\mathbb Q_2$.

First work with odd coefficients and put $H=G_F$.
Choose coherent roots of two.  Coherent odd roots of $\sqrt2$
can and will be chosen to square to them: the equation
$2a+mb=1$, for odd $m$, gives the unique compatible choice
at level $m$.  Thus
$\kappa_{\sqrt2}=\kappa_2/2$ literally.
All coherent odd roots of $\mu_8$ lie in $\mu_8\subset F$,
and their Kummer cocycles on $H$ vanish.

Here is the exact fourth-division construction of the cusp
path class.  Choose an $m$-division $v_m$ of the torus point
$4z_d=t_d$ using the roots just specified.  For integers
$a,b$ with $4a+mb=1$ set $u_m=a v_m+b z_d$ in the
generalized Jacobian.  Then
$m u_m=z_d$ and $4u_m=v_m$.
These points are compatible: the two required equations
determine $u_m$ uniquely, since multiplication by four is
an automorphism of the $m$-torsion.  Their proper images are
$p(u_m)=b[d-O]$, which are $H$-fixed four-torsion points;
thus their proper Kummer cocycle is literally zero.
They prove the exact
cocycle identity, in the fixed splitting (E103.16),
\[
       \kappa(z_d)(g)=v_d\kappa_2(g),\qquad
       v_d=\tfrac18 N_{\bullet d}\in W,
       \qquad p_*\kappa(z_d)=0.
 \tag{E105.13}
\]
In particular the vanishing compact coordinate follows from
the actual division construction.  The fact that $d-O$ is
killed by four alone would specify only its cohomology class.

The framed Abel--Jacobi comparison has the same explicit
vector-frame convention as (E24.3): for a function $f$ whose
divisor is supported at the cusps, its character is its value
on a divisor times the product of the frame ratios with the
specified divisor exponents.  Applied to $f_q$ this identifies
four times the $q$th torus coordinate; it is zero on $j(J)$
because $\operatorname{Hom}(J,\mathbb G_m)=0$.
Together with the proper Abel--Jacobi map these characters
detect $A_R$, by the actual inverse (E103.16).
Regularizing at both tangents yields (E105.10), with the
following exact vector-frame correction.  Put $\epsilon_d$
equal to $+1$ at $d$ and $-1$ at each other non-$O$ cusp.
For a unit with divisor $\sum m_q q$ its character on this
tuple is $(-1)^{m_d+m_O}$, since $\sum m_q=0$.
This is precisely the two endpoint sign from replacing
the local difference $s-t$ by $t-s$ in regularizing the
divisor frames, as in the framed comparison of E.37.
Thus the tangential class is $\kappa(z_d\epsilon_d)$;
the coherent odd divisions of $\epsilon_d$ are $H$-fixed.
This proves that (E105.13) is
the actual abelian path cocycle with those tangent and root
choices.  It also proves its sign without identifying an
unframed divisor with a tangential path.

Choose a full return path $p_d$ from the tangent at $d$ to
the tangent at $O$ lifting this abelian path.  The map from
the full path torsor to the abelian path torsor is onto;
starting with any path, multiply it by a lift of the required
abelian difference.  Let $d_d=p_d t_d^{\rm per}p_d^{-1}$,
where $t_d^{\rm per}$ is the positive local peripheral
generator.  The actual conjugator is
$u_d(g)=p_dg(p_d)^{-1}$, so
$u_d(g)g(d_d)u_d(g)^{-1}=d_d^{\chi(g)}$ and its abelian
coordinate is the positive cocycle in (E105.13).
These eleven peripherals extend to a free pro-$\ell$ basis.
Choose the six remaining abelian basis vectors in the
specific compact summand (E103.16) and lift them.  The
Frattini criterion gives a free basis, since the resulting
abelian basis change is invertible.  This is an actual
free-group basis morphism, retaining the peripheral paths.

\emph{The class-two action needed for the origin.}
For odd coefficients use BCH coordinates
$N_{R,2}=A\oplus\Lambda^2 A$ with multiplication
\[
 (x,z)(y,w)=(x+y,z+w+\tfrac12x\wedge y).
 \tag{E105.14}
\]
The conversion from ordered Hall coordinates is
$(a_i;b_{ij})\mapsto
(\sum a_ie_i,\sum b_{ij}e_i\wedge e_j+
\tfrac12\sum_{i<j}a_i a_j e_i\wedge e_j)$.
This gives its integral inverse and proves the claimed
class-two group identification.
Let $a(g)=\kappa_2(g)$ and define
$L:W\to\Lambda^2W$ by $L(e_d)=e_d\wedge v_d$.
Expanding the actual peripheral conjugation gives, on
$W\oplus\Lambda^2W$,
\[
 g(x,z)=(\chi_gx,\chi_g^2z+\chi_g a(g)L(x)).
 \tag{E105.15}
\]
Indeed $u_d^{-1}d_d^{\chi_g}u_d$ has central term
$\chi_g a(g)e_d\wedge v_d$.
The central coordinate of $u_d$ cancels.
Unknown action on the six other free generators does not
enter (E105.15); the formula proves that this displayed
subgroup is $H$-stable.

Choose the E.103 coset lifts to be the return paths from
the four conjugates $B_i$ of $O$, with the abelian
normalization just proved.  The averaged abelian origin
restricts to $H$ as $a(g)v$ relative to $s_H$, where
\[
 v=\tfrac14(v_{B_1}+v_{B_2}+v_{B_3})
   =\left(-\tfrac12,-\tfrac12,-\tfrac12,-\tfrac12,
          -\tfrac58,-\tfrac34,-\tfrac58,
          -\tfrac12,-\tfrac12,-\tfrac12,-\tfrac12\right).
 \tag{E105.16}
\]
To check the sign, its $i$th conjugate cusp section is
$u_i(h)s_H(h)$.  Thus the word $b_i(s_H(h))$ in
(E103.8) is $-u_i(h)$ in $A$, and (E103.9) adds their
average.  Both coordinate-unit rows kill $v$.
The tangential values at $O$ are $\zeta_8$ and the unit
leading coefficient one; their odd Kummer cocycles are
zero on $H$.  Hence $r_0|_H=0$.
A global $R(1)^2$ cocycle vanishing on $H$ is zero:
normality gives $h c(g)=c(g)$ for every $h\in H$, while
$R(1)^H=0$.  Consequently the toric correction in
(E103.10) vanishes for this normalization.

There is now a literal class-two $H$-section above this
abelian restriction.  Its geometric cocycle relative to
the original tangent section is
\[
        c_v(g)=\bigl(a(g)v,\ \tfrac12a(g)^2L(v)\bigr).
 \tag{E105.17}
\]
Substitution in (E105.14)--(E105.15) and
$a(gh)=a(g)+\chi_ga(h)$ proves
$c_v(gh)=c_v(g)g(c_v(h))$: the wedge of $v$ with
itself is zero and the cross term is
$\chi_ga(g)a(h)L(v)$.

Pull back $\mathcal E_{R,2}\to\mathcal E_R$ along the
\emph{fixed} abelian section $s_R$.
Its kernel is the abelian module $B_R$; the geometric
abelian part of an arithmetic lift acts trivially on
this central kernel.  Its $H$-section is (E105.17).
Apply the literal formula (E103.8)--(E103.9), now to
this abelian-kernel extension, using its four actual
cosets.  It produces a continuous global section of
the pullback and therefore the required $S_R$.
This is where division by four is used a second time.
It averages an extension with abelian kernel $B_R$,
not four sections of a nonabelian geometric extension.

Here is the promised comparison to every originally
chosen E.103 origin.  If $e_i'=a_i e_i$ are the aligned
coset lifts and $q=\tfrac14\sum_i a_i$, direct use of
(E103.8) gives
\[
 B'(e)=B(e)+g(4q)-4q,\qquad
 s_0'=\operatorname{Int}(q)s_0,\qquad
 s_R'=\operatorname{Int}(q-\iota r(q))s_R.
 \tag{E105.18}
\]
Lift $q-\iota r(q)$ to $N_{R,2}$ and conjugate the
constructed class-two section by its inverse.
Its abelian image is exactly the old $s_R$.
Changing unit roots by the cocycle coboundary $dk$
changes the toric-zero origin by
$\operatorname{Int}(\iota k)$, with the same lift rule.
These formulas retain all old choices, even when their
numerical path coordinates were not separately assigned.

\emph{The rational two-primary construction.}
For any $\ell$, the class-two Malcev group is the BCH
group of $A_{\mathbb Q_\ell}\oplus\Lambda^2A_{\mathbb Q_\ell}$.
The original $N_{\ell,2}$ maps into it by the displayed
Hall/BCH conversion.  Its arithmetic automorphisms extend
uniquely to this group.  The actual arithmetic pushout is
the quotient of $N_{\mathbb Q_\ell,2}\rtimes\mathcal E_{\ell,2}$
by the diagonal normal subgroup
$\{(n^{-1},n):n\in N_{\ell,2}\}$.
The diagonal subgroup is closed, and its intersection
with the first factor is trivial; this gives the specified
extension and the literal map from the original one.

Only the root normalization needs a new check at two.
On $H$, choose the coherent roots
$\zeta_{8\cdot2^n}$ of $\zeta_8$.  Then
$\kappa(\zeta_8^k)=k(\chi-1)/8$.
Choose compatible roots of $\sqrt2$ squaring to the
chosen roots of two.  If $u_n$ are any compatible
$2^n$-division points of $z_d$ and $v_n$ are the explicit
torus roots of $4z_d$, then
$b_d=(4u_n-v_n)_n\in T_2J_{C,D}$ is an actual Tate vector,
and
\[
 \kappa(z_d)=v_d\kappa_2+
        d\left(\tfrac14b_d+\tfrac1{32}K_{\bullet d}\right)
             \quad\text{in }A_{\mathbb Q_2}.
 \tag{E105.19}
\]
The compact coordinate in the first term is zero in
the fixed $j/4$ splitting.  Multiplying the return path
by the displayed vector subtracts precisely this
coboundary, since $p\mapsto m p$ changes
$p g(p)^{-1}$ by $m-gm$ in abelian coordinates.
For the actual tangential class add the endpoint correction
just specified: let $\sigma_d\in W$ have entry zero at $d$
and entry one at every other non-$O$ cusp.
Coherent roots of $-1$ give
$\kappa(\epsilon_d)=d(\sigma_d/2)$.
Hence the literal path shift is
$b_d/4+K_{\bullet d}/32+\sigma_d/2$.
The last term is retained separately because
$\epsilon_d^4=1$ and it is not detected by the matrix of
$4z_d$.  This is an actual rational class-two inner
transport; it need not be an integral pro-two path.

At $O$ the first coordinate unit has Kummer cocycle
$(\chi-1)/8$ and the second has zero.
Thus the aligned E.103 $r_0$ is exactly
$d(1/8,0)$ globally: its difference has zero restriction
to $H$ and $\mathbb Q_2(1)^H=0$.
Put $h_0=\iota(1/8,0)$ and transport both the tangent
section and the peripheral coordinates by
$\operatorname{Int}(h_0)$.
In these transported coordinates (E105.15)--(E105.17)
hold literally for the actual toric-zero origin.
All divisions by two and four exist, so the same
pulled-back central average proves the rational assertion.
Equation (E105.18) gives the inverse transport to every
specified old marking.  For odd rational coefficients
the earlier odd integral comparison followed by this
pushout gives the same construction.

\emph{Sufficiency and necessity for the complete toric image.}
The action relative to the new origin, restricted to $H$,
replaces $L$ by
\[
             L_v(w)=L(w)+v\wedge w.
 \tag{E105.20}
\]
The further central averaging does not change this action:
two lifts of the same abelian section differ centrally.
Let $w_A=\sum_i e_{A_i}$ and $w_C=\sum_i e_{C_i}$.
In $W$ the remaining orbit sum is $w_B=-w_A-w_C$.
The exact columns of (E105.12) give
\[
 L_v(w_A)=L_v(w_C)=0,\qquad
       \iota(u,v)=\tfrac14(uw_A+vw_B).
 \tag{E105.21}
\]
The program retains all $55\times11$ entries of $L,L_v$,
the two columns of $\iota$, their actual retraction, and
the literal zero matrix in (E105.21).

Suppose $\alpha\cup\beta=0$ and choose a continuous
cochain $k$ with $dk=\alpha\cup\beta$, using
$dk(g,h)=gk(h)-k(gh)+k(g)$.
On $H$, put $b=\iota(\alpha,\beta)$ and
$\omega=\iota(1,0)\wedge\iota(0,1)$.
The explicit central cochain
\[
           z(g)=-\bigl(k(g)+\alpha(g)\beta(g)/2\bigr)\omega
 \tag{E105.22}
\]
has differential equal to minus the BCH defect
$\tfrac12 b(g)\wedge g b(h)$.
Indeed $d(\alpha\beta)=-(\alpha\cup\beta+\beta\cup\alpha)$,
and (E105.21) removes the additional action term.
Thus $(b,z)S_R|_H$ is a genuine $H$-section above the
desired abelian input.  Pull back over that input and
apply the abelian-kernel four-coset formula once more.
This proves global sufficiency and gives a literal
construction of its lift.

For necessity use the actual morphism
$U\to\mathbb P^1-\{0,1,\infty\}$,
$t=-(x/z)^4$, $1-t=-(y/z)^4$.
It induces the class-two map and multiplication by four
on both Kummer coordinates.
The coherent roots of $-1$ give zero odd Kummer cocycle;
over $\mathbb Q_2$ their cocycle is the coboundary
$d(1/2)$, transported by the corresponding actual
inner change in the target BCH group.
The rational tangential paths at zero and one have
abelian Kummer class zero.  The standard based Galois
formula $x\mapsto x^\chi$, $y\mapsto f_g^{-1}y^\chi f_g$
has $f_g$ in the commutator subgroup, so its class-two
action is pure cyclotomic.  The explicit target
obstruction is therefore $16\alpha\cup\beta$; equivalently
this follows by multiplying its rank-two BCH law.
See also the actual formula and proof in
\cite[Proposition7, \S12.4.5]{E_Wickelgren}.
Since sixteen is invertible in every coefficient ring
in this theorem, a class-two section forces
$\alpha\cup\beta=0$.  This proves both image inclusions.

More generally the actual defect of a $W$-valued cocycle
$b$ on $H$, relative to $S_R$, is
\[
 \Omega_b(g,h)=\chi_g\kappa_2(g)L_v(b(h))
                     +\tfrac12 b(g)\wedge\chi_g b(h).
 \tag{E105.23}
\]
This specifies the action and cup terms for the other
toric directions, without replacing the handle action
by the associated-graded representation.

\emph{The continuous cup group, for every odd prime.}
We give the source calculation and its support maps,
since taking only finite or local cup tests would lose
part of (E105.5).
Put $G_S=G_{\mathbb Q,\{\ell,\infty\}}$ and
$T_n=\mathbb Z/\ell^n(2)$.
For $\ell=3$, the exact proof (E.32g) gives
$H^2(G_S,T_n)=H^3(G_S,T_n)=0$.
For $\ell\ge5$ let $K=\mathbb Q(\mu_\ell)$ and
$\Delta=\operatorname{Gal}(K/\mathbb Q)$, with character
$\omega$ on $\mu_\ell$.
Restriction and the $\Delta$-invariants functor are exact
here because $|\Delta|=\ell-1$ is invertible.
The literal coefficient identity
$\mathbb F_\ell(2)=\mu_\ell\otimes\mathbb F_\ell(1)$
and the Kummer sequence show that
$H^1(G_S,\mathbb F_\ell(2))$ is the $\omega^{-1}$ part
of the $S$-unit and $\ell$-torsion class groups of $K$.
The unique prime over $\ell$ is principal, so inverting it
does not change the class group.

The odd-character part of the $S$-unit group modulo
$\ell$ comes only from $\mu_\ell$, of character $\omega$.
For completeness, an $S$-unit $u$ has $u/\bar u$ integral
at every finite prime (the unique removed prime is fixed
by conjugation), and all its conjugates have absolute
value one.  It is a root of unity.  Hence conjugation is
trivial on the torsion-free unit quotient, and the stated
odd part follows.  Since $\omega^{-1}\ne\omega$ for
$\ell\ge5$, that $S$-unit component vanishes.
Ribet's original Main Theorem(1.1), with $k=2$, gives
the vanishing of the class component because $B_2=1/6$;
the used direction is Herbrand's theorem
\cite[p.151]{E105_Ribet1976}.
The theorem is stated for the class group modulo $\ell$;
its vanishing is equivalent to the vanishing of the
$\ell$-torsion component, since the $\Delta$ idempotents
act on the finite $\ell$-primary class group and Nakayama
applies to each summand.
We have proved $H^1(G_S,\mathbb F_\ell(2))=0$;
$H^0=0$ as well.  The real factor in the finite global
Euler characteristic is one, because this is an even
character.  Milne I.5.1 gives $H^2=0$.
Devissage and the odd-coefficient real-place formula
I.4.10(c), with finiteness I.4.15, now give
\[
 H^2(G_S,T_n)=H^3(G_S,T_n)=0,
                      \qquad H^1(G_S,T_n)\text{ finite}.
 \tag{E105.24}
\]

Apply the finite localization triangle on
$\operatorname{Spec}\mathbb Z[1/\ell]$ and remove a finite
set of primes $q\ne\ell$.
The support complex at $q$ is
$R\Gamma(\mathbb F_q,\mathbb Z/\ell^n(1))[-1]$:
tame inertia is $\mathbb Z_\ell(1)$, its wild subgroup has
prime-to-$\ell$ cohomology zero, and purity gives this
specified twist and shift.
First take the filtered direct limit over the removed
sets.  Finite continuous cochains factor through finite
extensions, so the generic-field term is exactly
$R\Gamma(G_{\mathbb Q},T_n)$ and the support term is
a direct sum.  Then take the derived inverse limit in
$n$, using \cite[Proposition1.6, Theorem2.2]{E_JannsenContinuous}.
The two terms in each residue complex are
\[
 [\,\mathbb Z/\ell^n\xrightarrow{q-1}\mathbb Z/\ell^n\,]
 \quad\text{in degrees }0,1,
 \tag{E105.25}
\]
with arithmetic Frobenius.  The direct sums have
surjective coefficient transitions in both terms.
Their derived inverse limit is the literal complex
$[\mathcal R\xrightarrow{d}\mathcal R]$, where
$\mathcal R$ consists of the $\mathbb Z_\ell$ sequences
tending to zero, indexed by $q\ne\ell$, and
$d(z)_q=(q-1)z_q$.
Equation (E105.24) and the finite degree-one groups kill
the adjacent base terms.  We obtain the actual module
isomorphism
\[
 H^2_{\rm cts}(G_{\mathbb Q},\mathbb Z_\ell(2))
       \simeq\mathcal R/d\mathcal R.
 \tag{E105.26}
\]
There is no closure in this denominator.
The localization, purity and arithmetic-\'etale maps used
here are Milne II.0.1(f), II.1.7(b), II.2.9; these are
the exact same maps retained in E.32.

To identify the cup at chain level, choose at $q$ tame
inertia $\tau$ and arithmetic Frobenius $\phi$ with
$\phi\tau\phi^{-1}=\tau^q$,
$\kappa_q(\tau)=1$, $\kappa_q(\phi)=0$.
The two Kummer cocycles on inertia are proportional.
Their alternating BCH cup therefore has a canonical
central-zero inertia splitting.  Write
$a_q,b_q$ for the two inertia values and $A'_q,B'_q$
for the Frobenius values.  In the rank-two BCH pullback,
with central inertia value $w$, its tame relation has
central defect
\[
             q\bigl((q-1)w+A'_qb_q-B'_qa_q\bigr).
 \tag{E105.27}
\]
The central Frobenius value cancels.
The support map divides this defect by $q$, since
$\operatorname{Hom}(\mathbb Z_\ell(1),T)=T(-1)$
compares its value at $\tau^q$ with its value at $\tau$.
Thus the cup maps to $[A'_qb_q-B'_qa_q]_q$ in (E105.26).
These inertia splittings are unramified outside a finite
set at each coefficient level.  They give a representative
in the direct sums before taking their derived limit.

Choose the place at $q$ to send the fixed
$\ell^{e_q}$th root of unity to $\eta_q$.
Finite-field Kummer theory then gives
$\kappa_p(\phi)\equiv r_{q,p}\pmod{q-1}$.
Write their difference as $(q-1)t_{q,p}$, with
$t_{q,p}\in\mathbb Z_\ell$.
The sums $\sum a_pt_{q,p}$ and $\sum b_pt_{q,p}$
converge and are uniformly bounded.
Replacing $A'_q,B'_q$ by $A_q,B_q$ changes the
representative by $d$ of a restricted family: every
term contains $a_q$ or $b_q$.
Changing global cocycles by
$\alpha\mapsto\alpha+du$, $\beta\mapsto\beta+dv$
adds $(q-1)(ub_q-va_q)$, again the image of a
restricted family.  Replacing $\phi$ by $\tau^r\phi$
leaves $A'_qb_q-B'_qa_q$ unchanged.
Changing $\eta_q$ multiplies the coordinate by a unit
and adds the preceding bounded-root correction.
These formulas prove the claimed independence with
the actual coboundary sign.
The unique possible inertia correction is (E105.4).
It belongs to $\mathcal R$ exactly under (E105.5),
proving the odd integral criterion.

\emph{Rational coefficients, with the same support topology.}
For every prime $\ell$, including two,
Hain--Matsumoto Theorem B.1 gives
$H^2_{\rm cts}(G_S,\mathbb Q_\ell(2))=0$;
their Proposition B.2 and Lemma B.3 prove the finite
coefficient comparison, and the proof on pp.38--39
includes the case $\ell=2$ \cite{E_HainMatsumoto}.
The degree-three base term also vanishes rationally:
Milne I.4.10(c) identifies finite-coefficient degrees
at least three with the real-place groups, which are
zero for odd $\ell$ and killed by two for $\ell=2$.
The finite-$S$ inverse systems are finite, so no
additional first derived limit occurs there.

For the unrestricted source one must retain a separate
cochain comparison.  For every degree $i$, compactness
of $G_{\mathbb Q}^i$ gives the isomorphism
\[
 C^i_{\rm cts}(G_{\mathbb Q},\mathbb Z_\ell(2))
           \otimes_{\mathbb Z_\ell}\mathbb Q_\ell
       \xrightarrow{\sim}
 C^i_{\rm cts}(G_{\mathbb Q},\mathbb Q_\ell(2)).
 \tag{E105.28}
\]
Every continuous cochain on the right has compact,
hence bounded, image, so multiplication by one power
of $\ell$ puts all its values in the integral lattice.
The same maps commute with differentials.
Apply this scalar extension to the finite-level
localization construction, after the derived inverse
limit.  The support complex becomes
$[c_0(\mathbb Q_\ell)\xrightarrow{d}c_0(\mathbb Q_\ell)]$,
and the rational base terms just computed vanish.
This proves the rational version of (E105.26).
The chain calculation (E105.27) remains literal;
it proves (E105.6).  A family tending to zero in
$\mathbb Q_\ell$ is bounded, so the condition retains
both a common denominator and convergence.  For the
product of odd rings, continuous cochains, their
differentials and products are componentwise; the
open-ideal topology in E.103 is exactly the resulting
product of the uniform primewise topologies.

For rational-number inputs the expression
$A_qb_q-B_qa_q$ is the discrete logarithm of (E105.7)
in its $\ell$-primary group.  This is a direct
finite-field calculation with the units obtained by
removing the $q$-valuation.  The sign has order two
and therefore vanishes for odd $\ell$.
There is finite support, so integrality is the entire
odd integral test and the rational test always holds.
This also agrees with the actual Milnor-symbol map
and tame-symbol computation in
\cite[Proposition26 and Remark28]{E_Wickelgren}; the
present proof does not infer an inverse-limit injection
from that finite-symbol statement.

To exhibit why convergence is necessary, use the
actual increasing prime sequence constructed in E.32:
$q_j\equiv1\pmod{3^j}$ and three is not a cube modulo
$q_j$.  Put $r_j=r_{q_j,3}$, which is a $3$-adic unit,
and set
\[
 a_3=1,\quad a_p=0\ (p\ne3),\qquad
 b_{q_j}=(q_j-1)/r_j,\quad b_p=0\text{ otherwise}.
 \tag{E105.29}
\]
This is a genuine restricted pair because
$v_3(b_{q_j})\ge j$.
Its support representative is $q_j-1$ at every $q_j$,
and its unique correction is $-1$ there.
At each finite coefficient level a finitely supported
primitive is obtained by putting $-1$ only where
$v_3(q_j-1)$ is smaller than that level.
Every local reduction is also zero.  The correction
fails to tend to zero even after multiplication by
any nonzero rational $3$-adic scalar.
Thus this pair has no integral or rational class-two
quartic lift.  Its abelian preimage exists by E.103,
so the obstruction has the asserted actual source.

\emph{Central fibres and the retained further target.}
The E.103 splitting identifies $A=W\oplus V^3$,
where $V=T_RE_+$ with its actual Weil pairing.
Exterior square gives $\Lambda^2W\oplus(W\otimes V^3)
\oplus\Lambda^2(V^3)$.
Each individual $\Lambda^2V_i$ is $R(1)$.
For each pair $i<j$, the two identical copies of $V$
identify $V_i\otimes V_j$ with
$\operatorname{Sym}^2V\oplus\Lambda^2V$ using
the projectors $(1\pm\mathrm{swap})/2$.
These are integral for odd $R$ and valid in every
rational target.  In a symplectic basis $(e,f)$,
the inverse column for the cross determinant is
$(e_i\wedge f_j-f_i\wedge e_j)/2$; contraction sends
it to one.  These six contraction/section maps and
their inverses are retained in the certificate.
This proves (E105.8) with actual coefficient maps.

There are no invariant vectors in $A$.
The toric part has none after restriction to $H$,
where it is $R(1)^{11}$.
For the proper part, the isogeny in E.103 and the
finite group $E_+(\mathbb Q)$ show that its integral
Tate invariants vanish; rational invariants vanish
after multiplying by a common denominator.
Consequently the stabilizer of an abelian section
under geometric conjugation is zero.
The exact central-extension sequence then makes
every nonempty fibre a torsor under the displayed
$H^1(\mathbb Q,B_R)$, with the stated continuous
cocycle action and its inverse given by the ratio
of two lifts.  To compare classes over an abelian
conjugacy class, first use the unique abelian conjugator
to align their actual abelian sections, and lift that
vector by the BCH coordinates (E105.14).  Its ambiguity
is central and is exactly the usual coboundary in $B_R$.
The alignment is continuous: for integral coefficients
the injective coboundary map has compact domain, and
for rational coefficients a finite collection of Galois
elements gives an injective linear matrix with a continuous
left inverse.  Thus the fibre maps also retain the stated
cochain-quotient topology.
The actual maps in (E105.8), and the explicit toric
defect (E105.23), specify what is evaluated here.
The remaining handle-action and cohomology calculations
for all abelian preimages are retained in E.23.
\end{proof}
